% !TeX root = ../nc-del-pezzo.tex
\section{中央多項式係数上の同じ ambient 代数における中心持上げ}
\label{sec:relative-center-lift}

ここでは，既に得た平坦な potential 代数の中心を持ち上げる．
別に選んだ星積の中心を，代数の同定なしに利用することはしない．
Etingof--Ginzburg の Proposition~8.1.1(i) は，形式 potential の中心持上げに
孤立特異点条件を課さず，その証明も Rees 変数を中央係数とする相対形式を用いる
\cite[\S8.1--8.2]{EG2010}．
以下では中央変数を任意個許し，同じ代数への同定，内部次数，
および有限パラメータへの移行を明記する．
使用する unimodularity 定理は複素数上の滑らかな affine 多様体についての
Dolgushev の定理であり，その未確認の相対版を仮定しない．

\subsection{形式中心と内部次数}

内部次数零の係数環と，正次数の中央変数を
\[
 R_0=\mathbb C[c_1,\ldots,c_N,\rho(c)^{-1}],\qquad
 R=R_0[s_1,\ldots,s_m],\qquad \deg s_j=v_j>0
\]
とする．$\rho=1$ を許し，$N,m$ は零でもよい．
$h$ の内部次数は零とし，活性変数の重みを
\[
 (w_x,w_y,w_z)=(1,1,2)\quad\text{または}\quad(1,2,3),
 \qquad d=w_x+w_y+w_z=2w_z
\]
とする．$Q=R[x,y,z]$ と置き，
\begin{equation}\label{eq:relative-center-classical-potential}
 F=\frac12z^2+U(s,x,y)\in Q_d,\qquad
 \Phi=xyz-yxz-h\operatorname{Sym}(F)
\end{equation}
を考える．$\operatorname{Sym}$ は活性変数の完全対称化であり，
$R$ に関して線型に延長する．
potential は巡回語として読み，その関係式は
\begin{equation}\label{eq:relative-center-potential-relations}
 [x,y]=hz,\qquad
 [y,z]=h\partial_x\operatorname{Sym}(U),\qquad
 [z,x]=h\partial_y\operatorname{Sym}(U)
\end{equation}
である．右辺の微分は巡回微分であり，可換化すると通常の偏微分となる．

\begin{lemma}[同じ相対 potential 代数の規格化形式中心]
\label{lem:relative-formal-center-lift}
$\widehat A$ を\eqref{eq:relative-center-potential-relations}の
$h$ 進完備な代数とする．次を仮定する．
\begin{enumerate}
 \item $\widehat A$ は $h$ 進分離完備かつ $h$-torsion-free であり，
       与えられた生成元による還元は $\widehat A/h\widehat A=Q$ である．
 \item $R[[h]]$ に相対的な対角は perfect であり，連続な完備包絡代数
       $\widehat A^{\mathrm e}_{R[[h]]}$ に関して
       \[
        \underline{R\operatorname{Hom}}_{\widehat A^{\mathrm e}_{R[[h]]}}
          (\widehat A,\widehat A^{\mathrm e}_{R[[h]]})
          \simeq\widehat A(d)[-3]
       \]
       という無捩れな相対 Calabi--Yau 同型を持つ．
\end{enumerate}
このとき，同じ $\widehat A$ の中に内部次数 $d$ の中心元
\[
 \widehat\Omega\in Z(\widehat A)_d,\qquad
 \widehat\Omega\bmod h=F
\]
が存在する．完全対称化による形式 PBW 座標で，その $z^2$ 係数を
ちょうど $1/2$ とすることができる．
各 $h$ 係数は $s_1,\ldots,s_m$ の多項式であり，中央変数の有理関数への
係数拡大は行わない．$U$ の最高次部分の判別式や $F$ の滑らかさは仮定しない．
\end{lemma}

\begin{proof}
\emph{形式 PBW 座標と相対 formality．}
完全対称化による連続な $R[[h]]$-線型写像
\[
 \sigma:Q[[h]]\longrightarrow\widehat A
\]
は還元後に恒等写像となる．全射性は $h$ 係数を順次持ち上げることで従い，
単射性は $h$-torsion-free 性と $h$ 進分離性から従う．
従って $\sigma$ は内部次数を保つ同型である．
ここでは平坦性を中心の存在から導いておらず，仮定1を先に使用している．
この座標で乗法を $Q[[h]]$ へ移すと，$R$ に関して線型で
内部次数零の形式的結合変形 $\mu_h$ を得る．

活性三変数の Kontsevich 写像を $R$-線型に延長して用いる
\cite[\S4, \S6]{Kontsevich}．
そのグラフは $x,y,z$ のみを微分し，係数 $c,s$ を微分しない．
相対多重微分作用素から全相対 Hochschild cochain への包含は
多項式 HKR の擬同型である．さらにこの擬同型と formality は
内部次数零の部分に制限しても擬同型である．
実際，相対 bar 分解と多項式 Koszul 分解の比較およびホモトピーを
次数を保つものに取ると，その比較は各内部次数で成立する．
多重微分作用素上では，線型な重み付き $\mathbb G_m$ 作用に関する
Kontsevich の同変性により全ての構造写像が内部次数を保存する．
従って $h$ 進 Maurer--Cartan 対応を内部次数零の複体に適用できる．
単位を保つ正規化 cochain を使えば，
\begin{equation}\label{eq:relative-center-formality-identification}
 \pi_h\in h\bigwedge^2_Q\operatorname{Der}_R(Q)[[h]],\qquad
 [\pi_h,\pi_h]=0,\qquad\deg\pi_h=0
\end{equation}
と，内部次数を保ち $R$ を固定する gauge 同型
\[
 T=1+O(h):(Q[[h]],\star_{\pi_h})
       \xrightarrow{\ \sim\ }(Q[[h]],\mu_h)
       \xrightarrow[\sigma]{\ \sim\ }\widehat A
\]
を得る．以下の $T$ はこの合成を表す．
一階交換子の規約を $[a,b]_{\star}=h\{a,b\}+O(h^2)$ と取ると，
\eqref{eq:relative-center-potential-relations}により
\[
 \pi_h=h\pi_F+O(h^2),\qquad
 \pi_F=F_x\partial_y\wedge\partial_z
        +F_y\partial_z\wedge\partial_x
        +F_z\partial_x\wedge\partial_y .
\]
以上は形式的な同型類の対応だけでなく，中心元を元の代数へ戻す
具体的な存在同型 $T$ を与える．

\emph{相対 CY から複素数上の絶対 CY へ．}
$\ell=N+m$ と置き，完備包絡代数を
\[
 B=\widehat A\widehat\otimes_{\mathbb C[[h]]}\widehat A^{\mathrm{op}},
 \qquad
 B_R=\widehat A\widehat\otimes_{R[[h]]}\widehat A^{\mathrm{op}}
\]
とする．$B_R$ は $B$ を中央差
\[
 c_i\otimes1-1\otimes c_i\quad(1\le i\le N),\qquad
 s_j\otimes1-1\otimes s_j\quad(1\le j\le m)
\]
で割ったものである．これらは中央正則列をなす．
実際，$h$ による還元後には局所化した多項式環の独立な中央差であり，
途中の各商も $h$-torsion-free となることを，形式 PBW 座標と
$h$ 係数による帰納法で確認できる．
この正則列の Koszul 分解により $B_R$ は $B$ 上 perfect であり，
内部次数のずれを省略すれば
\[
 R\operatorname{Hom}_B(B_R,B)\simeq B_R[-\ell]
\]
である．Koszul 最上外積は列挙した中央差の順序で自明化する．
導来 change of rings と仮定2から
\[
 \begin{split}
 R\operatorname{Hom}_B(\widehat A,B)
 &\simeq R\operatorname{Hom}_{B_R}
       (\widehat A,R\operatorname{Hom}_B(B_R,B))\\
 &\simeq\widehat A[-\ell-3]
 \end{split}
\]
を得る．対角の perfect 性もこの合成から従う．
内部次数を戻したずれは $d+\sum_jv_j$ であり，
左右の作用に自己同型の捻りはない．

\emph{絶対 unimodularity と体積形式の重み零成分．}
$X=\Spec Q$ は滑らかな複素 affine 多様体であり，
標準体積形式
\[
 \omega=dc_1\wedge\cdots\wedge dc_N
       \wedge ds_1\wedge\cdots\wedge ds_m
       \wedge dx\wedge dy\wedge dz
\]
はその正準束を自明化する．
\eqref{eq:relative-center-formality-identification}の $\pi_h$ を
$X$ 上の絶対二重ベクトルとも読む．
全変数についての Kontsevich グラフ公式は，この $\pi_h$ 上では
相対グラフ公式と一致する．各辺の微分添字は二重ベクトルの成分から来るため，
$\partial_{c_i}$ または $\partial_{s_j}$ を含む項が一つも生じないからである．
従って $\pi_h$ は，先ほどの絶対 CY 代数と同型な星積の
絶対 Kontsevich 類でもある．

Dolgushev の Theorem~3 は，正準束が自明な滑らかな複素 affine 多様体の
一変数形式量子化について，無捩れな Van den Bergh 双対化 bimodule と
その形式 Poisson 類の unimodularity が同値であることを述べる
\cite[Theorem~3]{DolgushevModular2009}．
同論文の Definition~1 と Proposition~1 により，この条件は
可逆元 $f\in Q[[h]]^\times$ が存在して
\begin{equation}\label{eq:relative-center-invariant-volume}
 L_{\pi_h}(f\omega)=0
\end{equation}
となることと同値である．今の場合は直前で全ての仮定を満たしている．

$Q$ の内部次数は非負であり，$Q_0=R_0$ である．
$f$ の各 $h$ 係数を内部次数零へ射影したものを
$f_0\in R_0[[h]]$ とする．$f\bmod h$ が $Q$ の単元であることから
$f_0\bmod h\in R_0^\times$ であり，$f_0$ は可逆である．
$\pi_h$ は内部次数零なので，\eqref{eq:relative-center-invariant-volume}の
内部次数 $\deg\omega=d+\sum_jv_j$ 成分は
\[
 0=L_{\pi_h}(f_0\omega)=f_0L_{\pi_h}\omega
\]
となる．最後の等号では $\pi_h$ が $R$ に相対的であることを使った．
従って $L_{\pi_h}\omega=0$ であり，活性成分以外を持たないことから
\[
 L_{\pi_h}(dx\wedge dy\wedge dz)=0
 \quad\text{in }\Omega^\bullet_{Q/R}[[h]]
\]
を得る．この重み射影によって，未知の体積形式を
標準の相対体積形式へ取り替えられる．

\emph{相対 Poincar\'e ホモトピーと Casimir．}
相対一形式
\[
 \alpha_h=\iota_{\pi_h}(dx\wedge dy\wedge dz)
     \in h\Omega^1_{Q/R}[[h]]
\]
は閉じており，内部次数 $d$ を持つ．
活性三変数の多項式 de Rham 複体において，通常の活性次数
（$x,y,z,dx,dy,dz$ をそれぞれ次数1とする）を使う．
その正次数 $r$ の成分上で，Euler 場
$E=x\partial_x+y\partial_y+z\partial_z$ による
$r^{-1}\iota_E$ は Poincar\'e ホモトピーを与える．
一形式の各単項式では $r>0$ なので，必要な割算は正整数だけである．
このホモトピーは中央係数を固定し，重み付き内部次数も保存する．
従って
\[
 \psi_h\in hQ_d[[h]],\qquad d_{Q/R}\psi_h=\alpha_h
\]
を取れる．$\alpha_h/h\bmod h=d_{Q/R}F$ なので，
$\psi_h$ に $hR_d$ の元を加えることにより
$\psi_h/h\bmod h=F$ とできる．
活性次元が3であるため，$\pi_h$ と $d_{Q/R}\psi_h$ の双対表示から
$[\pi_h,\psi_h]=0$ となり，$h$-torsion-free 性により
$\psi_h/h$ も $\pi_h$ の Casimir である．

\emph{中心へ戻す操作と規格化．}
$\pi_h$ で twist した Kontsevich の cochain 写像は，
Poisson cohomology の次数零を星積の Hochschild cohomology の次数零，
すなわちその中心へ送る．その $h$ による還元は HKR の恒等写像である．
従って Casimir $\psi_h/h$ の像を同型 $T$ で戻すと，
同じ $\widehat A$ の内部次数 $d$ の中心元で，還元が $F$ となるものを得る．
完全対称化座標での $z^2$ 係数を $a(h)$ とすれば，
$2w_z=d$ と全ての $s_j$ の正次数性により
\[
 a(h)\in R_0[[h]],\qquad a(h)\equiv\frac12\pmod h .
\]
$2a(h)$ は中央単元なので，得た元をこれで割れば所要の
$\widehat\Omega$ となる．
全ての段階で係数環は $R_0[s]$ のままであり，$s$ の分母は導入していない．
\end{proof}

\begin{remark}[平坦性・CY 性の入力と判別式]
前補題は，任意の potential 代数の平坦性を中心持上げから証明するものではない．
本稿の適用では，定理\ref{thm:simultaneous-degree-one-ambient}と
定理\ref{thm:simultaneous-poisson-ambient}の平坦な同時 ambient 族を，
Poisson 方向を自明化した chart に引き戻す．
その表示は\eqref{eq:relative-center-potential-relations}であり，
次数別平坦性から仮定1を満たす．
仮定2は既存の相対 potential 対角複体から得られる．
形式段階だけなら，先に確立した平坦性の下でこの複体と各項を $h$ 進完備化し，
還元後の三変数 Koszul 分解から cycle を順次持ち上げても完全性が従う．
巡回 Hessian の自己双対性は係数環上の恒等式であるため，無捩れな CY 同型を与える．
これは補題\ref{lem:elliptic-top-relative-cy}，
\ref{lem:degree-one-elliptic-top-relative-cy}および
普遍低次項の補題\ref{lem:elliptic-universal-lower-pbw}，
\ref{lem:degree-one-elliptic-universal-lower-pbw}と同じ機構であり，
楕円判別式を使わない形式段階の論法である．
とくに活性最高次式に重根がある場合にも，上の入力が満たされる限り
中心持上げを適用できる．
\end{remark}

\subsection{有限中心方程式による代数化}

\begin{proposition}[規格化中心の有限局所化による代数化]
\label{prop:relative-finite-center-lift}
$g\in R_0[h]$ は $g(c,0)\in R_0^\times$ を満たすとし，
$S=R_0[h,g^{-1}]$ とする．
中央変数 $s_1,\ldots,s_m$ と活性変数 $x,y,z$ による
同じ関係式\eqref{eq:relative-center-potential-relations}の
次数付き $S$-代数を $A$ とする．
各 $A_q$ は有限射影 $S$-加群であり，$h$ 進完備化は
補題\ref{lem:relative-formal-center-lift}の仮定を満たすと仮定する．
このとき，$f\in S$，$f\equiv1\pmod h$ と
\[
 \Omega\in Z(A_f)_d,\qquad
 \Omega\bmod h=\operatorname{Sym}(F)
\]
が存在する．有限個の必要次数で完全対称化を座標に取ることができ，
その座標における $\Omega$ の $z^2$ 係数は $1/2$ となる．
中心性，次数，還元および規格化の恒等式は，任意の可換 $S_f$-代数への
基底変更後にも成立する．基底変更の平坦性や被約性は必要ない．
追加の局所化は係数 $c,h$ だけに依存し，中央変数 $s$ は全ての値を許す．
\end{proposition}

\begin{proof}
$s$ も正次数の代数生成元として扱い，基底環 $S$ に含めないことが重要である．
各内部次数の可換多項式空間は有限自由 $S$-加群であり，
完全対称化はそれから $A_q$ への写像を与える．
必要な有限個の次数
\[
 q=d,\quad d+w_x,\quad d+w_y,\quad d+w_z
\]
でこの写像は $h=0$ 上の同型である．
有限射影加群間の写像が同型でない locus は閉集合であり，$V(h)$ と交わらない．
従って $1\bmod h$ となる一つの元を可逆化すれば，これらを同時に同型にできる．
実際，その閉集合の定義イデアル $I$ に対して $I+(h)=S$ なので，
$I$ 内に $1\bmod h$ となる元がある．
この局所化を済ませ，なお係数環を $S$ と書く．
$z^2$ 係数を読む $S$-線型写像を
$\lambda:A_d\to S$ とする．

$A/hA$ は可換であり，各次数の射影性により $A$ は $h$-torsion-free だから，
\[
 \begin{gathered}
 M=A_d,\qquad N'=A_{d+w_x}\oplus A_{d+w_y}\oplus A_{d+w_z},\\
 D:M\longrightarrow N',\qquad
 D(a)=h^{-1}([a,x],[a,y],[a,z])
 \end{gathered}
\]
は正則な $S$-線型写像である．ここで $h^{-1}$ は $hA$ 内での一意な割算を表し，
係数環では $h$ を可逆化していない．
$F_{\mathrm{sym}}=\operatorname{Sym}(F)\in A_d$ と書く．
一階交換子が Jacobian 括弧であり $F$ がその Casimir であることから
$D(F_{\mathrm{sym}})\in hN'$ である．

求める元を $\Omega=F_{\mathrm{sym}}+h\eta$ と置くと，
中心性と規格化は有限線型方程式
\begin{equation}\label{eq:relative-center-finite-linear-system}
 D\eta=-h^{-1}D(F_{\mathrm{sym}}),\qquad\lambda(\eta)=0
\end{equation}
になる．$s_j$ は既に中央なので，これら以外の交換子方程式は不要である．
写像と右辺を
\[
 L=(D,\lambda):M\longrightarrow N'\oplus S,\qquad
 b=\bigl(-h^{-1}D(F_{\mathrm{sym}}),0\bigr)
\]
と置く．前補題の規格化形式中心から
$\widehat\eta=(\widehat\Omega-F_{\mathrm{sym}})/h$ を作ると，
これは\eqref{eq:relative-center-finite-linear-system}の $h$ 進形式解である．

有限 $S$-加群 $C=\operatorname{coker}L$ における $b$ の像を $\bar b$ とする．
任意の素イデアル $\mathfrak p\supset(h)$ に対し，形式解は
$S_{\mathfrak p}$ の $h$ 進完備化へ写る．
Noether 局所環 $S_{\mathfrak p}$ のこの完備化は忠実平坦なので，
$\bar b_{\mathfrak p}=0$ となる．
従って有限巡回加群 $S\bar b$ の台は $V(h)$ と交わらず，
\[
 \operatorname{Ann}_S(\bar b)+(h)=S
\]
である．よって消滅イデアルの中に $f_1\equiv1\pmod h$ を取れ，
$S_{f_1}$ 上では $b$ は $L$ の像に入る．
その有限線型解 $\eta$ により所要の $\Omega$ を得る．
最初の座標のための局所化と合わせ，分母を払えば，元の $S$ の
$f\equiv1\pmod h$ による一回の主開局所化にまとめられる．

得られた中心性と規格化は有限表示代数内の実際の恒等式である．
従って任意の基底変更後にも保たれる．
この証明は形式級数への $h=1$ の代入ではなく，
有限加群の一つの余核類を消すことで有限パラメータの元を得るものである．
\end{proof}

\begin{remark}[この命題が与える範囲]
上の結果は，全内部次数4または6の中心元を，元の ambient 族の
零切断全体を含む有限主開集合上で与える．
局所的に選んだ Poisson 方向を変えたときの貼合せ，中心商の普遍正則性，
Noether 性，相対 tails および変形の完全性は，この命題の結論には含めない．
また，退化した古典 potential の Poisson 中心全体が $F$ で生成されるとは
主張していない．証明が使うのは一つの規格化 Casimir の存在だけである．
\end{remark}
